\documentclass[10pt]{article}
\usepackage[utf8]{inputenc}
\usepackage[T1]{fontenc}
\usepackage{amsmath}
\usepackage{amsfonts}
\usepackage{amssymb}
\usepackage[version=4]{mhchem}
\usepackage{stmaryrd}
\usepackage{bbold}

\begin{document}
L13

L13: Gauss' Lemma continued...

Proof (of Gauss' Lemma): Let $r_{1}, r_{2}, \ldots, r_{n}$ be\\
the lenst non-negative residues of the integers $a, 2 a, \cdots,\left(\frac{p-1}{2}\right) a$ that are greater than $P / 2$ and let $S_{1}, S_{2}, \ldots, S_{m}$ be the least non-negative residues of these integers that are less than $P / 2$ (no least positive residue is equal to $P / 2$ ). Note that no $r_{i}$ or $s_{j}$ is equal to zero, since $p$ does not divide any $a, 2 a, \ldots,\left(\frac{p-1}{2}\right) a$. Consider the $\frac{p-1}{2}$ integers given by

$$
p-r_{1}, p-r_{2}, \ldots, p-r_{n}, s_{1}, s_{2}, \ldots, s_{m} \text {. (1) }
$$

We wish to show that these integers are all the integers from 1 to $\frac{p-1}{2}$ inclusive (in some order).\\
Since each integer in (1) is between I and and $\frac{p-1}{2}$-inclusive, it suffices to prove\\
that no two of these integers are congruent modulo $p$. No two of the first n integers are congruent modulo $p$, otherwise $p-r_{i} \equiv p-r_{j}(\bmod p)$ for some $i \neq j$, and then $r_{i} \equiv r_{j}(\bmod p)$. Thus $k_{i} a \equiv k_{j} a(\bmod p)$ for some $k_{i}, k_{j} \in \mathbb{Z}$ with $k_{i} \neq k_{j}$ where $1 \leqslant k_{i} \leqslant \frac{p-1}{2}$ and $1 \leqslant k_{j} \leqslant \frac{p-1}{2}$. Since $p+a$, the inverse of a modulo $p$ exists, and thus $k_{i} \equiv k_{j}(\bmod p)$. Contradiction! Similarly, no two of the second $m$ integers are congruent modulo $P$. Finally, no one of the first $n$ integers can be congrnent modulo $p$ to one of the second $m$ integers, otherwise $p-r_{i} \equiv S_{j}(\bmod p)$ for some $i$ and $j$, and so $-k_{i} a \equiv k_{j} a(\bmod p)$ for some $k_{i}, k_{j} \in \mathbb{Z}$ with $k_{i} \neq k_{j}$ and $1 \leqslant k_{i}, k_{j} \leqslant \frac{p-1}{2}$. Applying the inverse of a\\
modulo p gives $-k_{i} \equiv k_{j}(\bmod P)$.

Contradiction 1\\
Thus the integers in (1) are the integers from 1 to $\frac{P-1}{2}$ inclusive (in Some order). Then\\
$\left(p-r_{1}\right)\left(p-r_{2}\right) \cdots\left(p-r_{n}\right) S_{1} S_{2} \cdots S_{m} \equiv\left(\frac{p-1}{2}\right)!(\bmod p)$.\\
Thus $(-1)^{n} r_{1} r_{2} \cdots r_{n} s_{1} s_{2} \cdots s_{m} \equiv\left(\frac{p-1}{2}\right)!(\bmod p)$.\\
By definition of the $r_{i}$ and $s_{j}$, we have that $(-1)^{n}(a)(2 a) \cdots\left(\frac{p-1}{2}\right) a \equiv\left(\frac{p-1}{2}\right)!(\bmod p)$,\\
or equivalently,

$$
(-1)^{n} a^{(p-1) / 2}\left(\frac{p-1}{2}\right)!\equiv\left(\frac{p-1}{2}\right)!(\bmod p)
$$

Since $\left(\left(\frac{P-1}{2}\right)!, P\right)=1$, the inverse of $\left(\frac{P-1}{2}\right)$ !\\
modulo $p$ exists, and thus we obtain

$$
\begin{aligned}
(-1)^{n} a^{(p-1) / 2} & \equiv 1(\bmod p), \\
\text { and so } \quad a^{(p-1) / 2} & \equiv(-1)^{n}(\bmod p) . \quad(2)
\end{aligned}
$$

Euler's criterion and (2) imply that\\
$\left(\frac{a}{p}\right) \equiv(-1)^{n}(\bmod p)$. Since $\left(\frac{a}{p}\right)= \pm 1$, we\\
have that $\left(\frac{a}{p}\right)=(-1)^{n}$, as desired.\\
Theorem 4.8: Let $p$ be an odd prime number.\\
Then $\left(\frac{2}{p}\right)=(-1)^{\left(p^{2}-1\right) / 8}=\left\{\begin{array}{cl}1 & \text { if } p \equiv 1,7(\bmod 8) \\ -1 & \text { if } p \equiv 3,5(\bmod 8) .\end{array}\right.$\\
Proof: By Gauss Lemma, we have that\\
$\left(\frac{\alpha}{p}\right)=(-1)^{n}$ where $n$ is the number of least\\
positive residues of the integers\\
$2,(2) 2, . . .\left(\frac{p-1}{2}\right) 2$ (3)\\
that are greater than $P / 2$. Let $k \in \mathbb{Z}$ with\\
$1 \leqslant k \leqslant \frac{P-1}{2}$. Then $k(2)<P / 2$ if and only if\\
$k<P / 4$. Thus only $\left[\frac{P}{4}\right]$ integers from (3)\\
are less than $P / 2$. Thus $\frac{P-1}{2}-\left[\frac{P}{4}\right]$ integers\\
In (3) are greater then $P / 2$, from which

$$
\left(\frac{2}{p}\right)=(-1)^{\frac{p-1}{2}-\left[\frac{p}{4}\right]} \text { by Ganss' Lemma. }
$$

It Suffices to prove that


\begin{equation*}
\frac{p-1}{2}-\left[\frac{p}{4}\right] \equiv \frac{p^{2}-1}{8}(\bmod 2) . \tag{4}
\end{equation*}


Case 1: $p \equiv 1(\bmod 8)$ : Then $p=8 k+1$ for Some $k \in \mathbb{Z}$, from which

$$
\begin{aligned}
\frac{p-1}{2}-\left[\frac{p}{4}\right] & =\frac{8 k+1-1}{2}-\left[\frac{8 k+1}{4}\right] \\
& =4 k-2 k=2 k \equiv 0(\bmod 2) .
\end{aligned}
$$

On the other hand,

$$
\frac{p^{2}-1}{8}=\frac{(8 k+1)^{2}-1}{8}=8 k^{2}+2 k \equiv 0(\bmod 2) .
$$

Thus (4) holds when $p \equiv 1(\bmod 8)$.\\
Establishing (4) in the cases $p \equiv 3,5,7(\bmod 8)$ follow by a similar computation and are left as an exercise. This establishes the first equality in the Theorem. The Second equality\\
in the Theorem follows from the first.

Example: We have $\left(\frac{2}{7}\right)=1$ by Theorem 4.8,\\
Since $7 \equiv 7(\bmod 8)$.

\begin{itemize}
  \item We have $\left(\frac{2}{53}\right)=-1$ by Theorem 4.8, since $53 \equiv 5(\bmod 8)$.\\
Recap: - We know how to compute $\left(\frac{-1}{p}\right)$ and $\left(\frac{2}{p}\right)$.
  \item We now need to understand $\left(\frac{q}{p}\right)$ for $p, q$ both odd primes and $p \neq q$.
\end{itemize}

Quadratic Reciprocity

\begin{itemize}
  \item The law of quadratic reciprocity relates\\
$\left(\frac{q}{p}\right)$ to $\left(\frac{p}{q}\right)$, which sometimes simplifies\\
Computations!\\
Theorem 4.9 (Law of Quadratic reciprocity)\\
Let $p$ and $q$ be distinct odd prime numbers.
\end{itemize}

Then\\
$\left(\frac{p}{q}\right)\left(\frac{q}{p}\right)=(-1)$\\
$\left(\frac{p-1}{2}\right) \cdot\left(\frac{q-1}{2}\right)$

$$
=\left\{\begin{array}{c}
1 \text { if } p \equiv 1(\bmod 4) \text { or } q \equiv 1 \bmod 4) \\
-1 \text { or both }) \\
\text { if } p \equiv q \equiv 3(\bmod 4)
\end{array} .\right.
$$

Remark: If $p \equiv 1(\bmod 4)$ or $q \equiv 1(\bmod 4)$, then $\left(\frac{p}{q}\right)\left(\frac{q}{p}\right)=1$ and so $\left(\frac{p}{q}\right)=\left(\frac{q}{p}\right)$.\\
If $p \equiv q \equiv 3(\bmod 4)$, then $\left(\frac{p}{q}\right)\left(\frac{q}{p}\right)=-1$, so that $\left(\frac{p}{q}\right)=-\left(\frac{q}{p}\right)$.\\
Example: Compute $\left(\frac{7}{53}\right)$. We have that

$$
\begin{aligned}
\left(\frac{7}{53}\right) & =\left(\frac{53}{7}\right) \quad(\text { by } T 4 \cdot 9 \text { and } 53 \equiv 1(\bmod 4) \\
& =\left(\frac{4}{7}\right) \\
& =7 \quad\left(\text { since } 4 \equiv 2^{2}(\bmod 7)\right) .
\end{aligned}
$$

Example: Compute $\left(\frac{-158}{101}\right)$.

Exercise:

$$
\begin{aligned}
\left(\frac{-158}{101}\right) & =\left(\frac{-1}{101}\right)\left(\frac{158}{101}\right) \\
& =\left(\frac{158}{101}\right) \\
& =\left(\frac{57}{101}\right) \\
& =\left(\frac{3}{101}\right)\left(\frac{19}{101}\right) \\
& =\left(\frac{101}{3}\right)\left(\frac{101}{19}\right) \\
& =\left(\frac{2}{3}\right)\left(\frac{6}{19}\right) \\
& =(-1)\left(\frac{6}{19}\right) \\
& =-\left(\frac{2}{19}\right)\left(\frac{3}{19}\right) \\
& =-(-1)\left(\frac{3}{19}\right) \\
& =\left(\frac{3}{19}\right) \\
& =-\left(\frac{19}{3}\right) \\
& =-\left(\frac{1}{3}\right)=-1 .
\end{aligned}
$$

Justify the\\
steps. steps.

Warning: It would not be correct\\
to write $\left(\frac{57}{101}\right)=\left(\frac{101}{57}\right)$ by Theorem 4.9. Why? 57 is not a prime number.


\end{document}